% bib_a 文献目录 Bibliography (书页 495-500 / p510-p515)

\chapter*{文献目录}
\markboth{文献目录}{}
\addcontentsline{toc}{chapter}{文献目录}

在正文中，我没有尝试为原始文献提供细致的引用。本着聚焦教学的原则，我在适当的地方给出了近期综述文章的参考文献。这里我列出若干书籍，它们或许可以作为你手中这本书的有用补充；这份清单无意做到包罗无遗，而是聚焦于那些仍在印行、且恰好为我所熟悉的书。

有两个网站是跟踪引力物理最新工作的宝贵资源。第一个是面向广义相对论与量子宇宙学的 ArXiv 电子预印本服务器：

\begin{center}
\texttt{http://arxiv.org/form/gr-qc/}
\end{center}

世界各地的研究者把他们最新的论文放在这里，随后便可以方便地下载。物理学其他领域也有类似的服务器。另一个网站是 Living Reviews in Relativity：

\begin{center}
\texttt{http://www.livingreviews.org/}
\end{center}

Living Reviews 是一份在线期刊，专门刊载引力物理各个领域的综述文章。对于任何有兴趣探究当前热门课题之最新工作的人来说，它都是一个绝佳的起点。

\subsection*{狭义相对论}

E. Taylor and J. Wheeler, \emph{Spacetime Physics} (Freeman, 1992)。一本非常好的狭义相对论入门书，花了极大的力气去化解这门学科似乎总要引发的种种“佯谬”。

A.P. French, \emph{Special Relativity} (W.W. Norton, 1968)。比 Taylor 和 Wheeler 的书稍显平淡，但却是通往狭义相对论的一条直截了当的入门之路。

\subsection*{广义相对论本科教材}

B.F. Schutz, \emph{A First Course in General Relativity} (Cambridge, 1985)。这是一本非常好的入门教材，真正努力在本科物理的常见课题与 GR 的语言和结果之间架起过渡的桥梁。

J.B. Hartle, \emph{Gravity: An Introduction to Einstein's General Relativity} (Addison-Wesley, 2002)。通过集中于弯曲时空的例子以及粒子在其中的行为，让 GR 的探索变得平易，只要可能就把物理放在形式之先。

E.F. Taylor and J.A. Wheeler, \emph{Exploring Black Holes: An Introduction to General Relativity} (Benjamin Cummings, 2002)。以黑洞为途径来介绍 GR 的物理原理。

\subsection*{广义相对论研究生教材}

R. Wald, \emph{General Relativity} (Chicago, 1984)。对许多高级课题——包括黑洞、整体结构和旋量——的透彻讨论。这是一部无比宝贵的参考书；如果你需要一个提得恰当的 GR 问题的正确答案，就该翻这本书。

C. Misner, K. Thorne and J. Wheeler, \emph{Gravitation} (Freeman, 1973)。这本书教育了至少两代引力物理研究者。内容全面而百科全书式，但文风时常独树一帜——你要么喜欢，要么不喜欢。

S. Weinberg, \emph{Gravitation and Cosmology} (Wiley, 1972)。在它所做的事情上是一本了不起的书，在天体物理、宇宙学和实验检验方面尤其出色。不过，它对这一题材采取了一种不寻常的非几何进路，而且不讨论黑洞。在一路推进那些令人印象深刻的大计算方面，Weinberg 比我们大多数人都强得多。

R. D'Inverno, \emph{Introducing Einstein's Relativity} (Oxford, 1992)。一部明智而清晰的广义相对论入门书，对现代 GR 课程所必需的主要课题给出了扎实的覆盖。

A.P. Lightman, W.H. Press, R.H. Price, and S.A. Teukolsky, \emph{Problem Book in Relativity and Gravitation} (Princeton, 1975)。收罗了 GR 各领域的大量习题，并附有演算完整的解答，这让教师们更难编出学生无法轻易查到答案的题目了。

\subsection*{广义相对论高级教程}

S. Hawking and G. Ellis, \emph{The Large-Scale Structure of Space-Time} (Cambridge, 1973)。一部强调整体技巧、微分拓扑与奇点定理的高级著作；一本经典。

F. de Felice and C. Clarke, \emph{Relativity on Curved Manifolds} (Cambridge, 1990)。取数学的进路，但对物理上可测量的量有着出色的强调。

R. Sachs and H. Wu, \emph{General Relativity for Mathematicians} (Springer-Verlag, 1977)。正如书名所言；不过，数学著作典型的干燥文风在这里被频繁插入的、带着鲜明个人判断的旁白激活了——既谈物理，也谈数学（还谈世界的状况）。

J. Stewart, \emph{Advanced General Relativity} (Cambridge, 2003)。对若干高级课题的简短而精到的介绍，尤其是旋量、渐近结构和特征初值问题。

\subsection*{数学背景}

B. Schutz, \emph{Geometrical Methods of Mathematical Physics} (Cambridge, 1980)。Schutz 的另一本好书，涵盖了那本 GR 书里略去的一些数学要点（但仍保持在非常容易入门的水平上）。其中包括关于李导数、微分形式的讨论，以及它们在 GR 之外的物理中的应用。

T. Frankel, \emph{The Geometry of Physics: An Introduction} (Cambridge, 2001)。一部内容丰富、可读性强的书，讲的是对物理真正有用的几何课题，包括流形、丛、曲率、李群与代数拓扑。

M. Nakahara, \emph{Geometry, Topology and Physics} (Institute of Physics, 2003)。一本平易近人的微分几何与拓扑入门书，重点放在物理学家感兴趣的课题上。

F.W. Warner, \emph{Foundations of Differentiable Manifolds and Lie Groups} (Springer-Verlag, 1983)。该领域的标准教材，既包括流形、张量场等基础课题，也包括更高级的内容。

\subsection*{专门课题}

J.D. Jackson, \emph{Classical Electrodynamics} (Wiley, 1999)。研究生水平电磁学的经典参考书。书中的习题在一代代研究生身上留下了难以磨灭的印记。

H. Goldstein et al., \emph{Classical Mechanics} (Prentice-Hall, 2002)。研究生水平力学的经典参考书。新修订版增加了更多关于非线性动力学的讨论。

V.I. Arnold, \emph{Mathematical Methods of Classical Mechanics} (Springer-Verlag, 1989)。对某些物理学家来说是一本令人生畏的书，但它从数学上老到的视角对经典力学所作的处理令人振奋。书里有大量很好的微分几何。

E.W. Kolb and M.S. Turner, \emph{The Early Universe} (Perseus, 1994)。已成为早期宇宙宇宙学的标准参考书，内容涉及暗物质、相变和暴胀。

A.R. Liddle and D. Lyth, \emph{Cosmological Inflation and Large-Scale Structure} (Cambridge, 2000)。聚焦于暴胀及其对大尺度结构的影响，对宇宙学微扰论作了细致的处理。

B.S. Ryden, \emph{Introduction to Cosmology} (Addison-Wesley, 2002)。一本非常新颖、注重物理的当代宇宙学课题入门书，面向高年级本科生或研究生新生。

S. Dodelson, \emph{Modern Cosmology} (Academic Press, 2003)。研究生水平的宇宙学入门书，着重于宇宙学微扰、大尺度结构和宇宙微波背景。

C.M. Will, \emph{Theory and Experiment in Gravitational Physics} (Cambridge, 1993)。一部有用的汇编，涵盖 GR 的各种替代理论以及对它们的实验约束，其中包括对参数化后牛顿形式的讨论。

S.L. Shapiro and S.A. Teukolsky, \emph{Black Holes, White Dwarfs and Neutron Stars: The Physics of Compact Objects} (Wiley, 1983)。一本自足的入门书，讲致密星与黑洞的物理与天体物理。

M.E. Peskin and D.V. Schroeder, \emph{An Introduction to Quantum Field Theory} (Westview Press, 1995)。已迅速成为量子场论的标准教材。

J. Polchinski, \emph{String Theory} (Cambridge, 1998)。现代弦论的标准两卷本入门书，包括对 D 膜与弦对偶的讨论。

C.V. Johnson, \emph{D-Branes} (Cambridge, 2003)。对被称为 D 膜的延展对象的详尽介绍——它们已成为弦论不可或缺的一部分；不需要预先具备弦论本身的知识。

E.E. Falco, P. Schneider, and J. Ehlers, \emph{Gravitational Lenses} (Springer Verlag, 1999)。对引力透镜理论与应用的一本透彻的入门书。

N.D. Birrell and P.C. Davies, \emph{Quantum Fields in Curved Spacetime} (Cambridge, 1984)。对于想要实用式入门弯曲时空量子场论（包括霍金效应）的人来说，这是标准读本。

R.M. Wald, \emph{Quantum Field Theory in Curved Spacetime and Black Hole Thermodynamics} (Chicago, 1994)。对弯曲时空中的量子场细致而在数学上严格的阐述；如果你真想知道真空态到底是什么，就看这里。

\subsection*{科普读物}

K.S. Thorne, \emph{Black Holes and Time Warps: Einstein's Outrageous Legacy} (W.W. Norton, 1994)。Thorne 是世界顶尖的引力物理研究者之一，涉猎各个方向；他既讲述了 GR 研究的历史，又介绍了非常前沿的研究课题。

R. Geroch, \emph{General Relativity from A to B} (Chicago, 1981)。对时空运行方式的一次真正漂亮的阐释。

B. Greene, \emph{The Elegant Universe: Superstrings, Hidden Dimensions, and the Quest for the Ultimate Theory} (W.W. Norton, 1999)。一本应时而又带着个人色彩的弦论物理入门书。它不怕讨论相当高深的概念，但始终面向普通读者；写得非常好。

A.H. Guth, \emph{The Inflationary Universe: The Quest for a New Theory of Cosmic Origins} (Perseus, 1998)。对整个现代宇宙学的透彻而清晰的介绍，重点在暴胀。

L. Smolin, \emph{Three Roads to Quantum Gravity} (Basic Books, 2002)。这“三条路”指的是弦论、圈量子引力，以及某种更深刻的东西；Smolin 是圈量子引力的拥趸，但这些讨论对所有人来说都应该是有趣的。

G. Kane, \emph{Supersymmetry: Unveiling the Ultimate Laws of Nature} (Perseus, 2001)。一本不错的超对称入门书。超对称是玻色子与费米子之间的一种假想对称性，也许很快就会进入粒子加速器所能企及的范围。

A. Einstein, H.A. Lorentz, H. Weyl, and H. Minkowski, \emph{The Principle of Relativity} (Dover, 1924)。其实根本不是一本“科普”书；它是一部狭义相对论与广义相对论的原始研究论文集，已译成英文。

A. Pais, \emph{Subtle Is the Lord: The Science and the Life of Albert Einstein} (Oxford, 1983)。一部 Einstein 的科学传记，方程俱全。
